\section{从垂直 Pipeline 到通用 Sequence Learning}

历史主线落到 NLP 时，最直观的变化不是模型参数增长，而是系统接口减少。2009 年的机器翻译需要 word alignment、phrase extraction、重排序、language-model rescoring 和大量独立特征；2013 年的研究社区也按 morphology、dialogue、discourse、sentiment、translation 等垂直方向组织。通用 sequence formulation 逐步把这些任务转化为共享表示学习问题。

\subsection{2009：神经网络只是复杂系统中的一个 Rescorer}

早期 statistical machine translation（SMT，统计机器翻译）通过多个可解释模块组合出完整系统。Neural probabilistic language model 当时常被放在 pipeline 尾部，为候选翻译重新打分，而不是端到端决定输出。这样的工程可以工作，但每个接口都要定义表示、搜索、特征和误差传递，知识被锁在专用组件里。

\lecturefigure{slide-02-sequence-models-2009.jpg}{2009 年机器翻译 pipeline 与 neural probabilistic language model}{官方录像恢复幻灯片 V002；00:05:09}

\begin{knowledgebox}{读图：寻找神经网络到底在哪里}
左图包含 alignment、phrase table、decoder 和 rescoring 等多个模块；右图的 neural language model 只负责其中一个分数。先数系统边界，再观察梯度能否跨边界端到端传播。图支持当时 neural component 已有价值，却没有控制整个任务，也解释了为何数据驱动 consolidation 会带来巨大工程收益。
\end{knowledgebox}

\teachervoice{Vaswani 回忆自己读博时的 MT 系统甚至比图中更复杂，并让全班寻找“神经网络在哪里”。答案是 Continuous State Language Model 主要用于 rescoring；今天看来，名称中特意强调 continuous，反而说明当时主流世界仍是离散 pipeline。}

\subsection{2013：任务与研究社区都被 Verticalize}

当模型不能共享足够多的表示与计算结构时，每个问题都需要自己的领域知识、特征和 evaluation culture。\term{verticalized system（垂直化系统）}指围绕某一任务建立专用数据、模块与指标的技术栈；它能深挖局部问题，却使跨任务 transfer 和共同基础设施变得困难。

\lecturefigure{slide-03-nlp-2013-verticalized.jpg}{NLP 2013 的研究轨道呈现高度垂直化}{官方录像恢复幻灯片 V003；00:07:03}

\begin{knowledgebox}{读图：会议目录也是技术组织图}
页面列出多个相对独立的 research track。先看同一语言现象被拆成多少任务，再问每个任务是否使用独立 feature、模型与数据。它不能说明专门研究没有价值，而是提示 consolidation 需要一种能共同处理 variable-length input、context 与 conditional generation 的通用计算接口。
\end{knowledgebox}

\subsection{Distributed Representation、Seq2Seq 与 Attention}

\term{distributed representation（分布式表示）}用稠密向量的多个维度共同编码概念，不再为每个词或规则维护独立符号特征；\term{contextual representation（上下文化表示）}进一步让同一 token 的向量随上下文改变。Seq2Seq 把 translation、summarization、dialogue、question answering 等任务统一为“输入序列到输出序列”，attention 则让 decoder 按内容读取 source states。

\lecturefigure{slide-04-general-approaches-2016.jpg}{从词向量到 Seq2Seq 与 Attention 的通用范式}{官方录像恢复幻灯片 V004；00:09:21}

\begin{knowledgebox}{术语消化：三次抽象升级}
\begin{tabularx}{\textwidth}{@{}L{0.20\textwidth}X X@{}}
\toprule
抽象 & 解决的问题 & 仍然留下的瓶颈 \\
\midrule
Distributed word vectors & 共享词语统计结构与语义相似性 & 同一词在不同上下文仍较固定 \\
Contextual sequence states & 让表示依赖前后文 & recurrent computation 仍按时间串行 \\
Seq2Seq + attention & 统一条件生成并按内容读取 source & encoder/decoder 主干仍常由 RNN 构成 \\
\bottomrule
\end{tabularx}
\end{knowledgebox}

\teachervoice{讲者把这一演化描述为“general paradigms started coming up”。关键不是 Word2Vec 的向量代数有多酷，而是研究者开始相信：只要能稳定学习 variable-length representation，许多原本独立的语言任务就能共享同一建模形式。}

\subsection{Variable-length Representation 是共同核心}

\term{sequence model（序列模型）}处理有顺序、长度可变的输入或输出。它必须把不同长度的 token sequence 转成可供预测使用的状态，同时保留内容、顺序和远距离依赖。下图用 recurrent encoder 表示经典路线：状态逐步更新，最终或逐位置表示被 decoder 与下游任务使用。

\lecturefigure{slide-05-variable-length-representations.jpg}{学习 variable-length representation 是 Seq2Seq 任务的共同基础}{官方录像恢复幻灯片 V005；00:10:33}

一个最小 recurrent update 为：
\[
h_t=f_\theta(h_{t-1},x_t),\qquad y_t\sim p_\theta(\cdot\mid h_t).
\]
其中 $x_t$ 是第 $t$ 个输入 token，$h_t$ 是汇总前缀 $x_{1:t}$ 的 hidden state，$f_\theta$ 是共享 recurrent transition，$y_t$ 是预测。递归让模型天然表达顺序，却使 $h_t$ 必须等待 $h_{t-1}$，因此并行深度随序列长度增长。

\begin{importantbox}{Consolidation 的第一步是统一问题表达}
在 Transformer 之前，Seq2Seq 已经完成重要抽象：translation、summarization 与 dialogue 都可以写成 conditional sequence modeling。Transformer 的贡献不是从零发明统一，而是把这套统一接口从 sequential backbone 中释放出来。
\end{importantbox}

\subsection{本章小结}

2009 的 MT pipeline 和 2013 的垂直研究轨道说明 specialization 曾经深入每个系统接口。Distributed/contextual representation、Seq2Seq 与 attention 先把任务形式统一；下一步的问题不是 RNN 是否有表达力，而是这种表达力能否在大规模数据和硬件上高效训练。

